\begin{exercise}*
  
  \begin{pycode}[2024-Quiz1-GrupoC-6]
A=Matrix([[1,2,3],[4,5,6],[7,8,9]]) & T((-1,1,3))
  \end{pycode}
  
  Suppose
  \begin{math}
    \Mat{A}= \py[2024-Quiz1-GrupoC-6]{latex(A)}
  \end{math}
  
  \begin{parts}
    
  \item \puntos{1} Encuentre una combinación lineal de las columnas de \Mat{A} que sea \Vect{0}.
    \begin{solution}[]
      Una posible respuesta es
      \begin{math}
        \;\Mat{A}\py[2024-Quiz1-GrupoC-6]{latex( (I(A.n) & T(Homogenea(A).pasos[1]))|0 )}=\Vect{0};\;
      \end{math}
      puesto que
      \begin{displaymath}
        \py[2024-Quiz1-GrupoC-6]{Homogenea(A).tex}
      \end{displaymath}
    \end{solution}
    
  \item \puntos{1} Explique por qué para cualquier matriz \Mat{B} de orden
    \begin{math}
      \py[2024-Quiz1-GrupoC-6]{A.n}\times\py[2024-Quiz1-GrupoC-6]{A.n},
    \end{math}
    el producto \Mat{AB} no tiene inversa.
    \begin{solution}[]
      Como las columnas de \Mat{A}
      (de orden $\py[2024-Quiz1-GrupoC-6]{A.n}\times\py[2024-Quiz1-GrupoC-6]{A.n}$)
      son linealmente independientes,
      generan un subespacio de dimensión menor que $\py[2024-Quiz1-GrupoC-6]{A.n}$.
      Y como las columnas de \Mat{AB} son combinaciones lineales de las de \Mat{A},
      también generan un subespacio de dimensión menor que $\py[2024-Quiz1-GrupoC-6]{A.n}$.
      Por tanto, \Matdim{AB}{\py[2024-Quiz1-GrupoC-6]{A.n}}{\py[2024-Quiz1-GrupoC-6]{A.n}} es singular.
      
      Alternativamente, como \Mat{A} es singular, existe un $\Vect{x}\ne\Vect{0}$ tal que \SELF{x}{A}{0}. Si \Mat{AB} fuera invertible, ello implicaría
      \begin{math}
        \Vect{x} = \VM{x}{AB}\InvMatp*{AB} =
        \Vect{0}\Mat{B}\InvMatp*{AB} = \Vect{0},
      \end{math}
      que contradice que $\Vect{x}\ne\Vect{0}$.
    \end{solution}
  \end{parts}
  \procedencia{MIT Course 18.06\quad Exam 1, Spring, 2021}
\end{exercise}

\begin{exercise}
  \puntos{0.5}
  Escriba un ejemplo de matriz de orden $5\times4$ con rango 3.
  \begin{solution}[]
    El ejemplo más sencillo es\;
    \begin{math}
      \footnotesize
      \begin{bmatrix}
        1&0&0&0\\
        0&1&0&0\\
        0&0&1&0\\
        0&0&0&0\\
        0&0&0&0\\
      \end{bmatrix}
    \end{math}
  \end{solution}
\end{exercise}

\begin{exercise}*
  
  \begin{pycode}[2025-Quiz1-GrupoH-1]
A = Matrix([
[ 1,  1,  1],
[ 1,  2,  3],
[ 1,  3,  5]])

B = A.span().cart  

from sympy import symbols

n = 3
b = Vector(symbols(f'b_1:{n+1}'))
condicion = -1 * A.amplia(-b) & T(SEL(A,b).pasos[1])|0|0

a = Vector([1,-2, 1])
c = Vector([2, 0,-2])
d = Vector([0,2,4,8])
e = Vector([1,-1,1,-1])
A_ampliada = A.amplia(c,1).amplia(a,1)
pasos_ampliada = A_ampliada.elim().pasos
  \end{pycode}
  
  Considere la matriz\;
  $\Mat{A}= \py[2025-Quiz1-GrupoH-1]{latex(A)}$
  
  \begin{parts}
  \item  \puntos{0.5}
    ¿Cuál es el rango de \Mat{A}?
    
    \begin{solution}[]
      Aplicando la eliminación vemos que
      $\rango{\Mat{A}}= \py[2025-Quiz1-GrupoH-1]{A.rango()}$\;
      {\small(emplear la matriz $[\Mat{A}|-\Vect{b}]$ nos ayudará a responder también al apartado b.)}
      \begin{displaymath}
        \footnotesize
        \py[2025-Quiz1-GrupoH-1]{Elim(A.amplia(-b,1)).tex}
      \end{displaymath}
    \end{solution}
    
  \item  \puntos{1} 
    Encuentre una matriz \Mat{B} tal que el espacio columna \cols{A} sea igual al espacio nulo \nulls{B}.
    
    \begin{solution}[]
      \eng
      De la última fila tras la reducción de $-\Vect{b}$ por eliminción,
      podemos ver que la condición para que un vector pertenezca al espacio columna de \Mat{A} es \;$\py[2025-Quiz1-GrupoH-1]{latex(condicion)}=0$.\;
      Así, para \;\fbox{$\Mat{B}=\py[2025-Quiz1-GrupoH-1]{latex(B)}$}\; tenemos
      \begin{displaymath}
        \cols{B}=
        \left\{
          \Vect{b}\in\R[3]
          \mid
          \py[2025-Quiz1-GrupoH-1]{latex(condicion)}=0 
        \right\}
        =
        \left\{
          \Vect{X}\in\R[3]
          \, \Big| \,
          \py[2025-Quiz1-GrupoH-1]{latex(B)}\Vect{x}=\Vect{0}
        \right\}
        =
        \nulls{B}.
      \end{displaymath}
    \end{solution}
    
  \item  \puntos{1}
    ¿Cuáles de los siguientes vectores pertennecen al espacio columna \cols{A}:
    \begin{displaymath}
      \py[2025-Quiz1-GrupoH-1]{latex(a)},\quad
      \py[2025-Quiz1-GrupoH-1]{latex(c)},\quad
      \py[2025-Quiz1-GrupoH-1]{latex(d)},\quad
      \py[2025-Quiz1-GrupoH-1]{latex(e)}\,?
    \end{displaymath}
    
    \begin{solution}[]
      \eng
      Los dos últimos vectores no pertenecen al espacio columna, pues son de \R[4].
      De la condición obtenida en la parte (b), tenemos que
      \begin{math}
        \py[2025-Quiz1-GrupoH-1]{latex(Vector(c,'h'))}
      \end{math}
      pertenece a \cols{A}, pues
      \begin{math}
        \Mat{B}
        \py[2025-Quiz1-GrupoH-1]{latex(Vector(c,'h'))}=
        \py[2025-Quiz1-GrupoH-1]{latex(B*c)}
      \end{math}
      pero
      \begin{math}
        \py[2025-Quiz1-GrupoH-1]{latex(Vector(a,'h'))}
      \end{math}
      no pertence puesto que.
      \begin{math}
        \Mat{B}
        \py[2025-Quiz1-GrupoH-1]{latex(Vector(a,'h'))}=
        \py[2025-Quiz1-GrupoH-1]{latex(B*a)}.
      \end{math}
      Otra forma de verlo es por eliminación:
      \begin{displaymath}
        \py[2025-Quiz1-GrupoH-1]{rprElim(A_ampliada, pasos_ampliada)}
      \end{displaymath}
    \end{solution}
  \end{parts}
  \procedencia{MIT Course 18.06\quad Quiz 1, March 5, 2007}
\end{exercise}

\begin{exercise}*
  
  \begin{pycode}[2025-Quiz1-GrupoH-2]
pasos = T([(-2,1,3), (1,3,2), (-1,3)], rpr='h')
E = I(3) & pasos
  \end{pycode}
  
  Sea \Matdim{A}{3}{3} tal que:
  \begin{math}
    \TrC[\py[2025-Quiz1-GrupoH-2]{latex(pasos)}]{\Mat{A}}=\Mat{I}.
  \end{math}
  
  \begin{parts}
  \item \puntos{1}
    Encuentre \InvMat{A}.
    
    \begin{solution}[]
      Las cuatro operaciones elementales por columnas en el proceso de eliminación Gauss-Jordan son\;
      \begin{math}
        \Mat{A}
        \TrCPE[\py[2025-Quiz1-GrupoH-2]{latex(pasos)}]{\Mat{I}}
        =\Mat{I};\,
      \end{math}
      es decir
      \begin{displaymath}
        \small
        \Mat{A}
        \py[2025-Quiz1-GrupoH-2]{latex(I(3) & pasos[0])}
        \py[2025-Quiz1-GrupoH-2]{latex(I(3) & pasos[1])}
        \py[2025-Quiz1-GrupoH-2]{latex(I(3) & pasos[2])}
        =
        \Mat{A}
        \py[2025-Quiz1-GrupoH-2]{latex(I(3) & pasos[0] & pasos[1])}
        \py[2025-Quiz1-GrupoH-2]{latex(I(3) & pasos[2])}
        =
        \Mat{A}
        \underbrace{\py[2025-Quiz1-GrupoH-2]{latex(E)}}_{\text{\InvMat{A}}}
        =\Mat{I}.
      \end{displaymath}
    \end{solution}

  \item \puntos{1}
    Encuentre \Mat{A}?
    
    \begin{solution}[]
      Aplicando sobre \Mat{I} la secuencia inversa de operaciones elementales
      (operaciones inversas en el orden inverso)
      podemos calcular la inversa de la matriz \InvMat{A} para obtener \Mat{A}:
      \begin{displaymath}
        \begin{bmatrix}
          \InvMat{A}\\\hline\Mat{I}
        \end{bmatrix}
        =
        \py[2025-Quiz1-GrupoH-2]{latex(E.apila(I(3),1))}
        \xrightarrow{\py[2025-Quiz1-GrupoH-2]{latex(T(pasos)**(-1))}}
        \py[2025-Quiz1-GrupoH-2]{latex(E.apila(I(3),1) & T(pasos**(-1)))}
        =
        \begin{bmatrix}
          \Mat{I}
          \\\hline
          \Mat{A}
        \end{bmatrix},
      \end{displaymath}
      O,también, calculando el siguiente producto:
      \begin{math}
        \;
        \TrCpE[\py[2025-Quiz1-GrupoH-2]{latex((pasos**-1)[0])}]{\Mat{I}}
        \TrCpE[\py[2025-Quiz1-GrupoH-2]{latex((pasos**-1)[1])}]{\Mat{I}}
        \TrCpE[\py[2025-Quiz1-GrupoH-2]{latex((pasos**-1)[2])}]{\Mat{I}}
        =\Mat{A};
      \end{math}
    \end{solution}
  \end{parts}
  \procedencia{Basado en MIT Course 18.06\quad Quiz 1, March 10, 1995}
\end{exercise}

\begin{exercise}*
  
  \begin{pycode}[2025-Quiz1-GrupoH-3]
L = ~Matrix([[1,2,2,1],[0,0,-4,-2],[0,0,0,0]])
U = Matrix([[1,0,3],[0,1,1],[0,0,1]])
A = L*U
c = sympy.symbols('c')
b = Vector([1,2,6,c])
eli = SEL(A,b,1)
sol = SEL(A,b,1,sust=[(c,3)])
  \end{pycode}
  
   Suponga que la matriz \Mat{A} es el siguiente producto de matrices: 
   \begin{displaymath}
     \Mat{A} = \MN{L}{U}=
     \py[2025-Quiz1-GrupoH-3]{latex(L)}
     \py[2025-Quiz1-GrupoH-3]{latex(U)};
     \quad
     \Parentesis{
       \text{\eng{así que}{so} }
       \Mat{A}\InvMatpE*{U}=\Mat{L}}.
   \end{displaymath}
 
   \begin{parts}
   \item \puntos{1}
     Obtenga una base del espacio fila de \Mat{A} y una base del espacio columna de \Mat{A}.
     
     \begin{solution}[]
       Como\;
       \begin{math}
         \Mat{A}=\Mat{L}\Mat{U}=
         \py[2025-Quiz1-GrupoH-3]{latex(L)}
         \py[2025-Quiz1-GrupoH-3]{latex(U)}
         =
         \py[2025-Quiz1-GrupoH-3]{latex(A)},\;
       \end{math}
       y como \Mat{L} tiene su primer pivote en la primera fila y el segundo pivote en la tercera fila; entonces
       \begin{displaymath}
         \text{\eng{Una base de}{A basis for} \cols{A}}:\;
         \py[2025-Quiz1-GrupoH-3]{latex((L|(1,2,)).sis())}.
         \qquad
         \text{\eng{Una base de}{A basis for} \cols{\MatT{A}}}:\;
         \py[2025-Quiz1-GrupoH-3]{latex(SubEspacio((~A).sis()).base)}.
       \end{displaymath}
     \end{solution}
     
   \item \puntos{1} 
     Describa explícitamente todas las soluciones del sistema \SEL{A}{x}{0}.
     
     \begin{solution}[]
       Como $\VectC{L}{3}=\elemR{\Mat{A}\InvMatpE*{U}}{3}$ es la única columna nula,\; $\big[\elemR{\InvMatpE*{U}}{3};\big]$ es una base de \nulls{A}.
       Como
       \begin{displaymath}
         \Mat{U}\InvMatpE*{U}=
         \py[2025-Quiz1-GrupoH-3]{latex(U)}
         \py[2025-Quiz1-GrupoH-3]{latex(InvMat(U))}
         =\Mat{I};
         \quad
         \text{\eng{entonces}{then}}
         \elemR{\InvMatpE*{U}}{3}=\py[2025-Quiz1-GrupoH-3]{latex(InvMat(U)|3)}.
       \end{displaymath}
       Así pues, el conjunto de soluciones del sistema \SEL{A}{x}{0} es\;
       \begin{math}
         \nulls{A}=
         \py[2025-Quiz1-GrupoH-3]{SubEspacio(A).EcParametricas()}.
       \end{math}
     \end{solution}
     
   \item 
     \puntos{1}
     Resuelva (si es posible, y dependiendo de $c$) el sistema
     \begin{math}
       \MV{A}{x}=\py[2025-Quiz1-GrupoH-3]{latex(Vector(b,'h'))}.
     \end{math}
     
     \begin{solution}[]
       $\py[2025-Quiz1-GrupoH-3]{eli.tex}$
       \medskip
       
       Si $c\neq3$ el sistema no tiene solución. 
       Cuando $c=3$ la solución completa es:
       \begin{displaymath}
         \py[2025-Quiz1-GrupoH-3]{latex(sol)}
       \end{displaymath}
     \end{solution}
     
   \end{parts}
   \procedencia{Basado en MIT Course 18.06\quad Quiz 1, March 10, 1995}
\end{exercise}
